\documentclass[11pt,a4paper]{article}
\usepackage[utf8]{inputenc}
\usepackage[T1]{fontenc}
\usepackage{amsmath,amssymb,amsthm}
\usepackage{listings}
\usepackage{geometry}
\usepackage{hyperref}
\geometry{margin=2.5cm}

\newtheorem{theorem}{Theorem}
\newtheorem{lemma}[theorem]{Lemma}
\newtheorem{observation}{Observation}

\lstset{
  language=C++,
  basicstyle=\ttfamily\small,
  keywordstyle=\bfseries,
  breaklines=true,
  numbers=left,
  numberstyle=\tiny,
  frame=single,
  tabsize=2
}

\title{IOI 2023 --- Overtaking}
\author{}
\date{}

\begin{document}
\maketitle

\section{Problem Statement}

$N$ buses and one reserve bus (bus $N$) travel a road with $M$ sorting
stations at positions $S[0] < S[1] < \cdots < S[M{-}1]$.  Bus $i$ has
speed $W[i]$ (time per unit distance) and departure time $T[i]$.  The
reserve bus has speed $X$ and variable departure time $Y$.

Between stations, each bus travels at its own speed.  At each station,
a bus's arrival time is the maximum of its own expected arrival and the
expected arrivals of all buses that arrived at the previous station
\emph{strictly before} it (slower buses ahead block faster ones).

For each query value of $Y$, determine the reserve bus's arrival time
at the last station.

\section{Solution}

\subsection{Key observation}

\begin{observation}
The reserve bus's presence does not affect the non-reserve buses'
mutual blocking order among themselves.  A slower non-reserve bus $i$
only blocks the reserve bus if $i$ arrived at the previous station
strictly before the reserve bus.
\end{observation}

This holds because the reserve bus can only block buses behind it (those
arriving later), but we only care about the reserve bus's own arrival
time, and the buses that block the reserve bus at each station are
exactly those non-reserve buses that arrived earlier.

\subsection{Precomputation}

Simulate the $N$ non-reserve buses through all $M$ stations (without
the reserve bus).  For each station $j$, store the arrival times sorted,
along with prefix maxima of their expected arrivals at station $j + 1$.

\subsection{Per-query processing}

For a given $Y$:
\begin{enumerate}
  \item Set $t_{\mathrm{res}} = Y$.
  \item For each segment from station $j$ to $j + 1$:
        \begin{enumerate}
          \item Compute the reserve bus's expected arrival:
                $e = t_{\mathrm{res}} + X \cdot (S[j{+}1] - S[j])$.
          \item Binary-search among non-reserve buses at station $j$ to find
                those arriving strictly before $t_{\mathrm{res}}$.
          \item The blocking value is the prefix maximum of their expected
                arrivals at station $j + 1$.
          \item $t_{\mathrm{res}} \gets \max(e, \text{blocking value})$.
        \end{enumerate}
  \item Return $t_{\mathrm{res}}$.
\end{enumerate}

\section{Implementation}

\begin{lstlisting}
#include <bits/stdc++.h>
using namespace std;
typedef long long ll;

int N, M;
ll X_speed;
vector<ll> S_pos;
// Per station j: sorted (arrival_time, expected_next) and prefix max
vector<vector<pair<ll,ll>>> sorted_st;
vector<vector<ll>> pfx_max;

void init(int L, int _N, vector<ll> T, vector<int> W,
          int _X, int _M, vector<int> S) {
    N = _N; M = _M; X_speed = _X;
    S_pos.resize(M);
    for (int i = 0; i < M; i++) S_pos[i] = S[i];

    // Simulate non-reserve buses
    vector<vector<ll>> arr(M, vector<ll>(N));
    for (int i = 0; i < N; i++) arr[0][i] = T[i];

    for (int j = 0; j + 1 < M; j++) {
        ll dist = S_pos[j+1] - S_pos[j];
        vector<pair<ll,ll>> av(N); // (arrival_j, expected_j+1)
        for (int i = 0; i < N; i++)
            av[i] = {arr[j][i], arr[j][i] + (ll)W[i] * dist};

        vector<int> ord(N);
        iota(ord.begin(), ord.end(), 0);
        sort(ord.begin(), ord.end(), [&](int a, int b){
            return av[a].first < av[b].first;
        });

        ll rmx = 0;
        for (int idx : ord) {
            ll e = av[idx].second;
            arr[j+1][idx] = max(e, rmx);
            rmx = max(rmx, e);
        }
    }

    // Build sorted station data
    sorted_st.resize(M);
    pfx_max.resize(M);
    for (int j = 0; j + 1 < M; j++) {
        ll dist = S_pos[j+1] - S_pos[j];
        vector<pair<ll,ll>> ae(N);
        for (int i = 0; i < N; i++)
            ae[i] = {arr[j][i], arr[j][i] + (ll)W[i] * dist};
        sort(ae.begin(), ae.end());
        sorted_st[j] = ae;
        pfx_max[j].resize(N);
        ll mx = 0;
        for (int i = 0; i < N; i++) {
            mx = max(mx, ae[i].second);
            pfx_max[j][i] = mx;
        }
    }
}

ll arrival_time(ll Y) {
    ll t = Y;
    for (int j = 0; j + 1 < M; j++) {
        ll dist = S_pos[j+1] - S_pos[j];
        ll expected = t + X_speed * dist;

        // Binary search: last bus at station j with arrival < t
        auto& ss = sorted_st[j];
        int lo = 0, hi = N - 1, pos = -1;
        while (lo <= hi) {
            int mid = (lo + hi) / 2;
            if (ss[mid].first < t) { pos = mid; lo = mid + 1; }
            else hi = mid - 1;
        }

        ll blocking = (pos >= 0) ? pfx_max[j][pos] : 0;
        t = max(expected, blocking);
    }
    return t;
}
\end{lstlisting}

\section{Complexity}

\begin{itemize}
  \item \textbf{Preprocessing:} $O(NM \log N)$ --- simulating $N$ buses
        through $M$ stations with sorting at each station.
  \item \textbf{Per query:} $O(M \log N)$ --- binary search at each of
        $M$ stations.
  \item \textbf{Space:} $O(NM)$.
\end{itemize}

\paragraph{Correctness.}  The reserve bus does not alter non-reserve
buses' arrival times in this formulation because we compute the
non-reserve simulation independently.  The only interaction is
one-directional: non-reserve buses block the reserve bus.  This is valid
because the problem's blocking rule only considers buses arriving
\emph{strictly before} the current bus, and adding the reserve bus does
not change the relative ordering among non-reserve buses.

\end{document}
